\chapter{The pure path}\label{chap:pure}

\section*{At a glance}
\begin{itemize}
\item \lead{Goal} The shipping code the final theorem is about: the dependency closure, a pure
  erasability oracle, a pure backend for the shipping traversal, and the entry point that routes
  in-scope inputs to it.
\item \lead{Main results} None: this chapter holds shipping definitions.
\item \lead{Status} \stPlanned{} every node: planned shipping entries S-18 to S-21 on branch
  \code{shipping}; the names are the planned names (namespace \code{Erasure}).
\item \lead{Lean files} \code{LeanToLambdaBox/Erasure.lean} and new shipping modules.
\item \lead{Reference} MetaRocq's \code{erase}, \code{is\_erasableb} and \code{erase\_global\_deps}
  (\code{erasure/theories/ErasureFunction.v}).
\end{itemize}

\section{Environment view and closure}\label{sec:pure-closure}

\begin{definition}[Errors of the pure path]
  \label{def:Erasure-EraseError}
  \lean{Erasure.EraseError}
  \planned
  \lead{In short} \code{outOfFragment} routes \code{\#erase} to the \code{Meta} path; the others
  (\code{nameCollision}, \code{fuel}, \code{failed}) are \hl{errors of \code{\#erase}}, never
  answers.
\end{definition}

\begin{definition}[Environment view]
  \label{def:Erasure-EnvView}
  \lean{Erasure.EnvView, Erasure.EnvView.ofEnvironment}
  \planned
  \lead{In short} What the pure path reads about the Lean environment: \code{find?},
  \code{isExtern}, \code{inlineAttr?}; the counterpart of MetaRocq's abstract environment.
  \begin{itemize}
  \item \code{ofEnvironment}: builds the view from the elaboration environment; \hl{trusted glue},
        as MetaRocq's quoting is trusted.
  \end{itemize}
\end{definition}

\begin{definition}[Lookup in a declaration list]
  \label{def:Erasure-findConst}
  \lean{Erasure.findConst}
  \planned
  \lead{In short} \code{findConst decls c}: the first declaration of \code{decls} named \code{c};
  MetaRocq's \code{lookup\_env}.
\end{definition}

\begin{definition}[Dependency closure]
  \label{def:Erasure-collectDeps}
  \lean{Erasure.collectDeps, Erasure.exprConsts, Erasure.declDeps, Erasure.closure,
    Erasure.findCollision, Erasure.collectFuel}
  \planned
  \uses{def:Erasure-EraseError, def:Erasure-EnvView}
  \lead{In short} \code{collectDeps view e}: the program's own environment, \hl{the closure of
  \code{e} through types, values and block members}, or the first out-of-fragment construct.
  \begin{itemize}
  \item \code{exprConsts}, \code{declDeps}: the names a term or declaration mentions, or
        \code{outOfFragment} at an inductive, constructor, recursor, quotient, literal,
        projection or metavariable;
  \item \code{closure}: depth-first, bounded by \code{collectFuel} $= 2^{32}$ steps; a view
        that answers with a declaration of another name is rejected;
  \item \code{findCollision}: two constants with the same $\lambda_\Box$ name (R-3) give
        \code{nameCollision}, after the fragment scan.
  \end{itemize}
  \lead{Reference} the environment \code{erase\_global\_deps} keeps, computed on the source side
  (DV-3).
\end{definition}

\section{The pure oracle}\label{sec:pure-oracle}

\begin{definition}[Universe instantiation]
  \label{def:Erasure-Pure-instLevels}
  \lean{Erasure.Pure.instLevels, Erasure.Pure.instLevel}
  \planned
  \lead{In short} Substitutes levels for level parameters in a term, \hl{without normalizing}
  levels, so that kernel reduction decides it.
  \lead{Reference} PCUIC \code{subst\_instance\_constr}.
\end{definition}

\begin{definition}[The pure erasability oracle]
  \label{def:Erasure-Pure-isErasable}
  \lean{Erasure.Pure.isErasable, Erasure.Pure.Ctx, Erasure.Pure.alwaysZero, Erasure.Pure.findLocal,
    Erasure.Pure.whnf, Erasure.Pure.inferType, Erasure.Pure.isArity, Erasure.oracleFuel}
  \planned
  \uses{def:Erasure-EraseError, def:Erasure-Pure-instLevels, def:Erasure-TravCtx}
  \lead{In short} \code{Pure.isErasable cx fuel ls e}: \code{e} is erasable if its type is an arity
  or its type's sort is evidently zero; every reduction uses \hl{the kernel's $\delta$}.
  \begin{itemize}
  \item \code{Ctx}: the program's declarations; \code{findLocal}: a traversal local by variable;
  \item \code{whnf}: head $\beta$, $\zeta$ (also of let-bound locals), $\delta$ of every
        definition, \code{mdata};
  \item \code{inferType}: infer-only retyping; binders the oracle enters are de Bruijn, free
        variables are the traversal's locals;
  \item \code{isArity} after \code{whnf}; \code{alwaysZero}: \code{zero}, \code{max} of two,
        \code{imax} with a right one;
  \item fuel exhaustion and ill-typed input are errors; \code{oracleFuel} $= 2^{20}$ per call.
  \end{itemize}
  \lead{Reference} \code{is\_erasableb} (\code{erasure/theories/ErasureFunction.v}); MetaCoq paper
  Section 7.2, Fig.~17; the procedure is ours (DV-18).
\end{definition}

\section{The pure backend and the entry point}\label{sec:pure-backend}

\begin{definition}[The pure backend]
  \label{def:Erasure-PureM}
  \lean{Erasure.PureM, Erasure.PureCtx, Erasure.PureState, Erasure.EraseT.runPure,
    Erasure.PureM.findConst?, Erasure.PureM.declInfo?, Erasure.PureM.unsafeRecBase?,
    Erasure.PureM.freshFVarId, Erasure.PureM.instantiate1, Erasure.PureM.isErasable,
    Erasure.PureM.casesInfo?, Erasure.PureM.ctorArity?, Erasure.PureM.prepare}
  \planned
  \uses{def:Erasure-EraseT, def:Erasure-findConst, def:Erasure-Pure-isErasable,
    def:Erasure-EnvView}
  \lead{In short} $\mathrm{PureM} = $ \code{ReaderT PureCtx (StateT PureState (Except
  EraseError))}: the traversal's backend \hl{without \code{Meta}}.
  \begin{itemize}
  \item \code{PureCtx}: the closure and the view; \code{PureState}: a counter for fresh
        variables;
  \item \code{findConst?} $=$ \code{declInfo?}: lookup in the closure, no \code{\_unsafe\_rec}
        redirect; \code{instantiate1}: a shipping copy of the pure \code{Expr.instantiate1'} of
        lean4lean's model, so that shipping code imports no lean4lean module;
  \item \code{isErasable}: the pure oracle at \code{oracleFuel}; \code{casesInfo?},
        \code{ctorArity?}: none;
  \item \code{prepare}: the identity (no \code{macro\_inline}, matcher inlining or \code{csimp}).
  \end{itemize}
\end{definition}

\begin{definition}[Recursion and remapping tests]
  \label{def:Erasure-isRecursiveDecl}
  \lean{Erasure.isRecursiveDecl, Erasure.nameOccurs, Erasure.axiomatized}
  \planned
  \uses{def:Erasure-EnvView, def:Erasure-ErasureConfig}
  \lead{In short} The traversal's tests at the pure backend, with exact names.
  \begin{itemize}
  \item \code{isRecursiveDecl}: a block member, or a declaration whose value mentions its own name
        (\code{nameOccurs}); such constants are compiled to \code{tFix} (DV-7);
  \item \code{axiomatized}: the configuration remaps the \code{@[extern]} declaration to a
        $\lambda_\Box$ axiom (DV-12).
  \end{itemize}
\end{definition}

\begin{definition}[Erasure on the pure path]
  \label{def:Erasure-erasePure}
  \lean{Erasure.erasePure}
  \planned
  \uses{def:Erasure-erase, def:Erasure-PureM, def:Erasure-isRecursiveDecl, def:ASTType}
  \lead{In short} \code{erasePure view cfg decls e}: \hl{the shipping traversal run with the pure
  backend} on the closure \code{decls}, assembling the program as \code{erase} does.
  \lead{Reference} MetaRocq \code{erase} and \code{erase\_global\_deps}, the pair of
  \code{erase\_correct}.
\end{definition}

\begin{definition}[The entry point]
  \label{def:Erasure-eraseEntry}
  \lean{Erasure.eraseEntry, Erasure.route, Erasure.Route}
  \planned
  \uses{def:Erasure-collectDeps, def:Erasure-erasePure, def:Erasure-erase}
  \lead{In short} What \code{\#erase} runs: \code{collectDeps}, then
  \begin{itemize}
  \item \code{ok decls}: \code{erasePure}; its errors are errors of \code{\#erase}, with
        \hl{no fallback to \code{Meta}};
  \item \code{outOfFragment}: the \code{CoreM} path, unchanged;
  \item any other error: an error of \code{\#erase}.
  \end{itemize}
\end{definition}

\section*{Caveats}
\begin{itemize}
\item \lead{Caveat} The pure path skips Lean's compiler passes, so its output may differ from the
  \code{Meta} path's on in-scope inputs; the register records each difference when S-21 lands.
\item \lead{Caveat} \code{EnvView.ofEnvironment} decides which constants are \code{@[extern]}, and
  so which constants the source semantics treats as stuck.
\end{itemize}
